%%%%%%%%%%%%%%%%%%%%%%%%%%%%%%%%%%%%%%%%%%%%%%%%%%%%%%%%%%%%%%%%%%%%%%%%%%%%
% CAPÍTULO 1: INTRODUCCIÓN Y MOTIVACIÓN
%%%%%%%%%%%%%%%%%%%%%%%%%%%%%%%%%%%%%%%%%%%%%%%%%%%%%%%%%%%%%%%%%%%%%%%%%%%%

\chapter{Introducción y motivación}
\label{ch:introduccion}

\section{Contexto}

La movilidad universitaria constituye un caso particular dentro de los sistemas de transporte urbano. Diariamente, un número elevado de estudiantes y personal se desplaza entre sus lugares de residencia y el campus universitario, generando patrones de desplazamiento con características propias: horarios relativamente concentrados, orígenes dispersos y destinos comunes, y una fuerte recurrencia semanal. En muchos casos, varios miembros de la comunidad universitaria realizan trayectos similares en franjas horarias próximas sin coordinación entre ellos. Como consecuencia, un mismo trayecto puede ser cubierto por varios vehículos de forma independiente, lo que reduce la ocupación media de los vehículos, incrementa la demanda de aparcamiento y aumenta el número total de desplazamientos asociados al campus.

\section{Problema}

El problema que motiva este Trabajo Fin de Máster (TFM) puede formularse de la siguiente manera: existen desplazamientos universitarios potencialmente compartibles que no se coordinan por falta de mecanismos específicos que permitan identificar de forma eficiente qué usuarios podrían compartir un trayecto. Aunque existen plataformas generales de coche compartido (\textit{carpooling}), estas no están adaptadas al contexto universitario ni incorporan mecanismos que permitan representar digitalmente el sistema de movilidad y evaluar su comportamiento antes de un despliegue real.

\section{Motivación}

La motivación del trabajo se apoya en cuatro ideas complementarias:
\begin{enumerate}
    \item \textbf{Sostenibilidad y eficiencia.} El coche compartido puede reducir el número de vehículos en circulación, mejorar la ocupación media y disminuir la presión sobre el aparcamiento del campus.
    \item \textbf{Contexto favorable.} La comunidad universitaria comparte horarios, destinos y patrones recurrentes, lo que facilita la formación de grupos de viaje relativamente estables.
    \item \textbf{Necesidad institucional.} El Plan Director Smart Campus de la Universidad de Almería \cite{UALSmartCampus2019} incluye entre sus actuaciones la adquisición de una solución para fomentar el uso del coche compartido (actuación n.º 43) y propone como indicadores el número de viajes en coche compartido y su evolución anual. El propio plan señala que la red de coche compartido disponible funciona de forma aislada de los sistemas universitarios y no permite controlar su efectividad, que es precisamente lo que aporta un Gemelo Digital con métricas de evaluación.
    \item \textbf{Oportunidad tecnológica.} Los Gemelos Digitales han demostrado su utilidad en sistemas de movilidad y transporte para monitorización, análisis y simulación de escenarios. Su aplicación específica al coche compartido universitario entre particulares (P2P, del inglés \textit{peer-to-peer}) permanece poco explorada.
\end{enumerate}

\section{Enfoque del trabajo}

El TFM une dos líneas de trabajo complementarias: la literatura sobre Gemelos Digitales aplicados a movilidad y transporte, y la literatura sobre coche compartido, \textit{ride-sharing} y movilidad compartida. El resultado es una propuesta que integra ambos elementos en un caso de uso concreto: una plataforma universitaria de coche compartido P2P con viajes previamente planificados, apoyada por un Gemelo Digital que representa el estado del sistema y permite simular su comportamiento.

La revisión sistemática de la literatura realizada confirma que la combinación de \textit{Digital Twin} y movilidad compartida ya ha sido estudiada en contextos como flotas centralizadas, vehículos autónomos o transporte bajo demanda. Sin embargo, su adaptación al modelo P2P universitario, donde el conductor ya tiene previsto realizar el desplazamiento y comparte plazas disponibles, no ha sido abordada de forma específica. Esta adaptación constituye el espacio que el presente TFM pretende cubrir.
